\documentclass[10pt,a4paper,fontset=none]{ctexart}
\setCJKmainfont{Noto Serif CJK SC}[BoldFont={Noto Sans CJK SC}]
\setCJKsansfont{Noto Sans CJK SC}
\setCJKmonofont{Noto Sans Mono CJK SC}
\usepackage[margin=22mm,headheight=14pt]{geometry}
\usepackage{amsmath,amssymb}
\usepackage{longtable,booktabs,array,calc}
\usepackage{graphicx}
\usepackage{xcolor}
\usepackage{xurl}
\usepackage{fancyhdr}
\usepackage{enumitem}
\usepackage{needspace}
\usepackage{fvextra}
\usepackage{hyperref}
\definecolor{reportblue}{HTML}{164B70}
\hypersetup{colorlinks=true,linkcolor=reportblue,urlcolor=reportblue,
  pdftitle={面向进迭时空 K1 的 Qwen3.5
本地推理优化},pdfauthor={riscv-accl project}}
\urlstyle{tt}
\ctexset{section/format=\Large\bfseries\raggedright,
  subsection/format=\large\bfseries\raggedright}
\setlength{\parindent}{0pt}
\setlength{\parskip}{4pt}
\setlength{\emergencystretch}{3em}
\setlength{\LTpre}{8pt}
\setlength{\LTpost}{8pt}
\setcounter{secnumdepth}{-1}
\setcounter{tocdepth}{1}
\renewcommand{\arraystretch}{1.10}
\setlist{nosep,leftmargin=1.5em}
\pagestyle{fancy}
\fancyhf{}
\fancyhead[L]{\small\textcolor{reportblue}{riscv-accl / SpaceMiT K1}}
\fancyhead[R]{\small 2026-10-16}
\fancyfoot[C]{\small\thepage}
\providecommand{\tightlist}{\setlength{\itemsep}{0pt}\setlength{\parskip}{0pt}}
\usepackage{color}
\usepackage{fancyvrb}
\newcommand{\VerbBar}{|}
\newcommand{\VERB}{\Verb[commandchars=\\\{\}]}
\DefineVerbatimEnvironment{Highlighting}{Verbatim}{commandchars=\\\{\}}
% Add ',fontsize=\small' for more characters per line
\newenvironment{Shaded}{}{}
\newcommand{\AlertTok}[1]{\textcolor[rgb]{1.00,0.00,0.00}{\textbf{#1}}}
\newcommand{\AnnotationTok}[1]{\textcolor[rgb]{0.38,0.63,0.69}{\textbf{\textit{#1}}}}
\newcommand{\AttributeTok}[1]{\textcolor[rgb]{0.49,0.56,0.16}{#1}}
\newcommand{\BaseNTok}[1]{\textcolor[rgb]{0.25,0.63,0.44}{#1}}
\newcommand{\BuiltInTok}[1]{#1}
\newcommand{\CharTok}[1]{\textcolor[rgb]{0.25,0.44,0.63}{#1}}
\newcommand{\CommentTok}[1]{\textcolor[rgb]{0.38,0.63,0.69}{\textit{#1}}}
\newcommand{\CommentVarTok}[1]{\textcolor[rgb]{0.38,0.63,0.69}{\textbf{\textit{#1}}}}
\newcommand{\ConstantTok}[1]{\textcolor[rgb]{0.53,0.00,0.00}{#1}}
\newcommand{\ControlFlowTok}[1]{\textcolor[rgb]{0.00,0.44,0.13}{\textbf{#1}}}
\newcommand{\DataTypeTok}[1]{\textcolor[rgb]{0.56,0.13,0.00}{#1}}
\newcommand{\DecValTok}[1]{\textcolor[rgb]{0.25,0.63,0.44}{#1}}
\newcommand{\DocumentationTok}[1]{\textcolor[rgb]{0.73,0.13,0.13}{\textit{#1}}}
\newcommand{\ErrorTok}[1]{\textcolor[rgb]{1.00,0.00,0.00}{\textbf{#1}}}
\newcommand{\ExtensionTok}[1]{#1}
\newcommand{\FloatTok}[1]{\textcolor[rgb]{0.25,0.63,0.44}{#1}}
\newcommand{\FunctionTok}[1]{\textcolor[rgb]{0.02,0.16,0.49}{#1}}
\newcommand{\ImportTok}[1]{#1}
\newcommand{\InformationTok}[1]{\textcolor[rgb]{0.38,0.63,0.69}{\textbf{\textit{#1}}}}
\newcommand{\KeywordTok}[1]{\textcolor[rgb]{0.00,0.44,0.13}{\textbf{#1}}}
\newcommand{\NormalTok}[1]{#1}
\newcommand{\OperatorTok}[1]{\textcolor[rgb]{0.40,0.40,0.40}{#1}}
\newcommand{\OtherTok}[1]{\textcolor[rgb]{0.00,0.44,0.13}{#1}}
\newcommand{\PreprocessorTok}[1]{\textcolor[rgb]{0.74,0.48,0.00}{#1}}
\newcommand{\RegionMarkerTok}[1]{#1}
\newcommand{\SpecialCharTok}[1]{\textcolor[rgb]{0.25,0.44,0.63}{#1}}
\newcommand{\SpecialStringTok}[1]{\textcolor[rgb]{0.73,0.40,0.53}{#1}}
\newcommand{\StringTok}[1]{\textcolor[rgb]{0.25,0.44,0.63}{#1}}
\newcommand{\VariableTok}[1]{\textcolor[rgb]{0.10,0.09,0.49}{#1}}
\newcommand{\VerbatimStringTok}[1]{\textcolor[rgb]{0.25,0.44,0.63}{#1}}
\newcommand{\WarningTok}[1]{\textcolor[rgb]{0.38,0.63,0.69}{\textbf{\textit{#1}}}}
\RecustomVerbatimEnvironment{Highlighting}{Verbatim}{breaklines=true,fontsize=\small,commandchars=\\\{\}}
\begin{document}
\begin{titlepage}
\vspace*{15mm}
{\large\textcolor{reportblue}{RISC-V 本地推理}}\par
\vspace{12mm}
{\Huge\bfseries 面向进迭时空 K1 的 Qwen3.5 本地推理优化\par}
\vspace{8mm}
{\Large 基础设施优化与性能评估\par}
\vspace{12mm}
{\large 性能证据截止：2026-10-16（Asia/Singapore）；后续验收更新：2026-10-22\par}
\vspace{14mm}
\begin{tabular}{@{}ll@{}}
\textbf{平台} & MUSE-Pi-Pro，K1／X60，标称 16 GB 内存 \\
\textbf{模型} & Qwen3.5 2B 和 4B \\
\textbf{实现} & SpaceMiT llama.cpp 分支 \\
\textbf{内容} & 中文技术报告 \\
\end{tabular}
\vfill
\textcolor{reportblue}{\rule{\textwidth}{1pt}}\par
本报告比较最终默认构建与厂商起点。证据附录见英文报告。\par
\vspace{6mm}
\href{https://github.com/MoyaMryia/riscv-accl}{github.com/MoyaMryia/riscv-accl}
\end{titlepage}
{\small\tableofcontents}
\clearpage
性能证据截止：2026-10-16，Asia/Singapore。后续验收记录更新：2026-10-22。

\hypertarget{ux6458ux8981}{%
\section{摘要}\label{ux6458ux8981}}

最终版完成一次请求的速度，2B 为起点的 1.24～3.59 倍，4B 为起点的
1.41～4.72 倍。 起点为厂商 \nolinkurl{5ad05d8} 加两项兼容修改。
输入越长，收益越大。 服务器内存基本不变，两臂之差小于 3 MiB。

\begin{longtable}[]{@{}lrrrrrr@{}}
\toprule
模型 & 输入词元 & 配对数 & 请求时长，s & 加速倍数 & 首词元时间减少 &
解码，词元/s\tabularnewline
\midrule
\endhead
2B & 512 & 3 & 43.031／34.714 & 1.24 倍 & 17.96\% &
3.387／4.292\tabularnewline
2B & 2048 & 3 & 136.554／97.169 & 1.41 倍 & 29.41\% &
3.108／4.186\tabularnewline
2B & 8192 & 1 & 875.359／426.221 & 2.05 倍 & 51.77\% &
2.239／3.600\tabularnewline
2B & 16384 & 2 & 2868.832／1057.226 & 2.71 倍 & 63.43\% &
1.659／2.914\tabularnewline
2B & 32768 & 1 & 10822.886／3014.282 & 3.59 倍 & 72.27\% &
1.051／2.141\tabularnewline
4B & 512 & 3 & 113.128／80.346 & 1.41 倍 & 21.72\% &
1.315／2.133\tabularnewline
4B & 2048 & 3 & 411.793／246.579 & 1.67 倍 & 39.77\% &
1.133／1.965\tabularnewline
4B & 8192 & 1 & 3323.894／1170.324 & 2.84 倍 & 65.17\% &
0.724／1.485\tabularnewline
4B & 16384 & 2 & 11388.985／3033.905 & 3.75 倍 & 73.56\% &
0.490／1.117\tabularnewline
4B & 32768 & 1 & 42294.220／8952.505 & 4.72 倍 & 78.93\% &
0.291／0.752\tabularnewline
\bottomrule
\end{longtable}

斜线前为起点均值，斜线后为最终版均值。
加速倍数为起点请求时长除以最终版请求时长。 两臂使用相同的 Q4\_0
权重、F16 缓存、输入和 65 个输出词元。 完整数据见第 5 节。

版本 \nolinkurl{k1-final-20261013-r3} 是两个模型共用的一份构建。
选定参数写在源码中，普通启动不需要性能参数或实验变量。
该构建通过原生检查、状态检查、自然代码回答和迁移启动检查。
默认配置是冻结搜索中的均衡选择。 样本数和待验收项见文末``适用范围''。

\hypertarget{ux8981ux6c42ux4e0eux72b6ux6001}{%
\section{1. 要求与状态}\label{ux8981ux6c42ux4e0eux72b6ux6001}}

\href{https://github.com/MoyaMryia/riscv-accl/blob/main/Requirements.md}{Requirements.md}
是唯一的现行验收依据。
它要求两个指定基座模型、本地推理、厂商分支、匹配对照提升和可复现部署材料。
它没有规定 8 GB 设备或 32k／64k 上下文。
完整要求对照见\href{https://github.com/MoyaMryia/riscv-accl/blob/main/spacemit/reports/2026-10-13-requirements-gap.md}{缺口报告}。

\begin{longtable}[]{@{}ll@{}}
\toprule
状态 & 内容\tabularnewline
\midrule
\endhead
已实现 &
一套完整源码、一份生产构建、固化默认值、依赖检查和统一启动器\tabularnewline
已验证 & 两个现用 GGUF 在 K1 上本地运行，不调用云端 API\tabularnewline
已验证 & 28 次普通默认性能请求、12 次容量扩展请求和 6
次完整自然代码回答\tabularnewline
已验证 & 原生转储逐位一致、64
条短状态记录和两个模型的迁移启动\tabularnewline
已验证 & 两个官方检查点的转换、量化和打包；产物与历史 GGUF
分别记录\tabularnewline
已验证 & 固定多轮 v2 完成 48 轮请求与独立复核；2B 15/24、4B 19/24
正确，用户接受测得质量\tabularnewline
已验证 & 自然长上下文在短、16k、32k 档位共 12
次回答全部正确\tabularnewline
待验收 & 长状态正确性和 8192／16384／32768 的重复稳定性\tabularnewline
本版未实现 & GPU 加速、KV 压缩和分页分配的内存收益\tabularnewline
\bottomrule
\end{longtable}

本项目没有对外比赛提交记录。

\hypertarget{ux5e73ux53f0ux4e0eux8f6fux4ef6ux8eabux4efd}{%
\section{2.
平台与软件身份}\label{ux5e73ux53f0ux4e0eux8f6fux4ef6ux8eabux4efd}}

测试板为 MUSE-Pi-Pro，有 8 个 K1／X60 CPU 核。 标称内存为 16 GB，Linux
报告约 15.528 GiB。 记录的 CPU 频率为 1.6 GHz。 测试在 tmux
中以普通用户执行，并持有板端共用锁。
本轮测试没有修改设备权限或系统设置。

\begin{longtable}[]{@{}ll@{}}
\toprule
\begin{minipage}[b]{0.47\columnwidth}\raggedright
项目\strut
\end{minipage} & \begin{minipage}[b]{0.47\columnwidth}\raggedright
身份\strut
\end{minipage}\tabularnewline
\midrule
\endhead
\begin{minipage}[t]{0.47\columnwidth}\raggedright
最终版本\strut
\end{minipage} & \begin{minipage}[t]{0.47\columnwidth}\raggedright
\nolinkurl{k1-final-20261013-r3}\strut
\end{minipage}\tabularnewline
\begin{minipage}[t]{0.47\columnwidth}\raggedright
厂商源码\strut
\end{minipage} & \begin{minipage}[t]{0.47\columnwidth}\raggedright
\texttt{a990751d\allowbreak{}55a4c54a\allowbreak{}cf2bb77d\allowbreak{}44282c20\allowbreak{}93652359}\strut
\end{minipage}\tabularnewline
\begin{minipage}[t]{0.47\columnwidth}\raggedright
完整源码文件数\strut
\end{minipage} & \begin{minipage}[t]{0.47\columnwidth}\raggedright
3271\strut
\end{minipage}\tabularnewline
\begin{minipage}[t]{0.47\columnwidth}\raggedright
最终补丁 SHA-256\strut
\end{minipage} & \begin{minipage}[t]{0.47\columnwidth}\raggedright
\texttt{b03d193f\allowbreak{}618d8303\allowbreak{}ce563c4f\allowbreak{}cd4e0372\allowbreak{}6a1f6f45\allowbreak{}2dcc10ec\allowbreak{}edb1a793\allowbreak{}8550477b}\strut
\end{minipage}\tabularnewline
\begin{minipage}[t]{0.47\columnwidth}\raggedright
服务器 SHA-256\strut
\end{minipage} & \begin{minipage}[t]{0.47\columnwidth}\raggedright
\texttt{f254f174\allowbreak{}919c0a97\allowbreak{}ef34366f\allowbreak{}116705f6\allowbreak{}d22574ca\allowbreak{}b4c2bfc3\allowbreak{}9928c0d3\allowbreak{}15b5bfbd}\strut
\end{minipage}\tabularnewline
\begin{minipage}[t]{0.47\columnwidth}\raggedright
配置记录 SHA-256\strut
\end{minipage} & \begin{minipage}[t]{0.47\columnwidth}\raggedright
\texttt{aeef14e2\allowbreak{}356dc9e3\allowbreak{}3b1fc6ab\allowbreak{}595545eb\allowbreak{}67af8189\allowbreak{}4f2c7cbd\allowbreak{}28294862\allowbreak{}8453bdb0}\strut
\end{minipage}\tabularnewline
\begin{minipage}[t]{0.47\columnwidth}\raggedright
编译器\strut
\end{minipage} & \begin{minipage}[t]{0.47\columnwidth}\raggedright
GCC 14.2.0，含厂商 IME 支持\strut
\end{minipage}\tabularnewline
\begin{minipage}[t]{0.47\columnwidth}\raggedright
运行库\strut
\end{minipage} & \begin{minipage}[t]{0.47\columnwidth}\raggedright
SPERT 0.6.2 和
\nolinkurl{libspine_tcm.so.3.0.1}，文件散列见构建收据\strut
\end{minipage}\tabularnewline
\begin{minipage}[t]{0.47\columnwidth}\raggedright
构建\strut
\end{minipage} & \begin{minipage}[t]{0.47\columnwidth}\raggedright
Release，IME1 开，OpenMP 关，Native 关，ZBA 开，repack 关，CPU
后端\strut
\end{minipage}\tabularnewline
\bottomrule
\end{longtable}

两个模型文件均为 Q4\_0 权重和 Q4\_0 嵌入。 文件名含
MTP，表示文件带有草稿张量；本版使用直接解码。
交付包不含模型权重和厂商运行库。

\begin{longtable}[]{@{}ll@{}}
\toprule
\begin{minipage}[b]{0.47\columnwidth}\raggedright
模型\strut
\end{minipage} & \begin{minipage}[b]{0.47\columnwidth}\raggedright
GGUF SHA-256\strut
\end{minipage}\tabularnewline
\midrule
\endhead
\begin{minipage}[t]{0.47\columnwidth}\raggedright
\nolinkurl{Qwen3.5-2B-MTP-Q4_0-embQ4_0-dv64k.gguf}\strut
\end{minipage} & \begin{minipage}[t]{0.47\columnwidth}\raggedright
\texttt{56f0ddd9\allowbreak{}0dfa0e24\allowbreak{}56af6cb4\allowbreak{}718ece41\allowbreak{}9678d21a\allowbreak{}75d65e27\allowbreak{}353319b5\allowbreak{}e431a2fa}\strut
\end{minipage}\tabularnewline
\begin{minipage}[t]{0.47\columnwidth}\raggedright
\nolinkurl{Qwen3.5-4B-MTP-Q4_0-embQ4_0-dv64k.gguf}\strut
\end{minipage} & \begin{minipage}[t]{0.47\columnwidth}\raggedright
\texttt{0adf6cce\allowbreak{}5df53921\allowbreak{}606033c3\allowbreak{}9a13fae2\allowbreak{}054543b7\allowbreak{}8af0e7ff\allowbreak{}98860e86\allowbreak{}f078ac0f}\strut
\end{minipage}\tabularnewline
\bottomrule
\end{longtable}

\href{https://github.com/MoyaMryia/riscv-accl/blob/main/spacemit/reports/2026-10-16-final-release-tcm.md}{最终版
TCM 跟踪}核对了两次短请求的设备访问。
两个服务器都加载了固定库，但访问硬件 TCM 返回 \nolinkurl{EACCES}。 13
份进程映射中没有 \nolinkurl{/dev/tcm}。 \nolinkurl{/dev/tcm_sync_mem}
不存在，共享屏障回退到堆内存。
这两次请求只用于跟踪设备访问，没有产生新的性能或质量结果。

\hypertarget{ux8d77ux70b9ux4e0eux517cux5bb9ux4feeux6539}{%
\section{3.
起点与兼容修改}\label{ux8d77ux70b9ux4e0eux517cux5bb9ux4feeux6539}}

起点为 SpaceMiT 分支的
\texttt{5ad05d87\allowbreak{}96cb094e\allowbreak{}e23a3a01\allowbreak{}52874f8a\allowbreak{}2b3060f8}。
原版 K1 构建纳入了一个 IME2 文件，构建失败。 第一项兼容修改只删除这一个
CMake 源文件列表项。 修改后的构建无法加载当前模型的缩减 MTP 草稿头。
第二项兼容修改读取 \nolinkurl{d2t}，并按存储宽度加载草稿头。
两臂都没有启用草稿解码。

修改只涉及 CMake 列表和 \nolinkurl{src/models/qwen35.cpp}
中的张量加载函数。 主输出头、计算图代码和全部计算内核保持原版字节。
兼容补丁 SHA-256 为
\texttt{2e8c40a4\allowbreak{}a0567873\allowbreak{}80e37750\allowbreak{}bab00577\allowbreak{}7580978d\allowbreak{}3d9a58d9\allowbreak{}01e911dc\allowbreak{}12fe01d7}。
两次失败和两项兼容修改的记录都保存在证据中。

两臂都使用 4 线程、批大小 32、上下文 12288、单槽和 F16 键值缓存。
两臂使用相同的编译器和运行库。
测得的收益是从原始软件到最终软件的整体收益。
历史权重量化的变化不计入收益。

\hypertarget{ux5747ux8861ux9ed8ux8ba4ux914dux7f6eux7684ux9009ux62e9}{%
\section{4.
均衡默认配置的选择}\label{ux5747ux8861ux9ed8ux8ba4ux914dux7f6eux7684ux9009ux62e9}}

冻结搜索使用一份已核验的 RC1 构建，共 104 次请求。 4
次预检在原版和候选构建上检查两个模型。
掩码阶段在两个模型、两种长度上测试 8 种组合，各两轮。 参数阶段测试 2／4
线程和 16／32／64／128 批大小。 4 次自然回答完成选型检查。 参数阶段保留
4 线程和批大小 32，没有触发额外复验或回退。

均衡分数对两个模型、两种长度、首词元时间和解码耗时赋相同权重。
首词元时间指客户端从发出请求到收到第一个输出词元的时间。 与最低分相差
3\% 以内的组合，按冻结规则视为无法区分。
在该区间内，规则优先选择启用候选较少的组合。 掩码 7 的观测分数最低，为
0.8904876472483813。 掩码 3 的分数为
0.915536301962804，处于该区间内，因此规则选择掩码 3。
选型分数相对同构建的 C0 对照计算，与第 5 节的真实起点对照分开。

\begin{longtable}[]{@{}ll@{}}
\toprule
编译默认值 & 取值\tabularnewline
\midrule
\endhead
工作线程与批线程 & 4／4\tabularnewline
批大小与微批大小 & 32／32\tabularnewline
上下文与槽位 & 12288／1\tabularnewline
缓存与注意力 & F16 键值缓存，闪存注意力开\tabularnewline
GPU 层数与解码 & 0，直接解码\tabularnewline
公共路径 & 宽 RVV 注意力和 IME route 3 开\tabularnewline
候选掩码 & 3：B1-r2 和 C1 开，C2-r2 关\tabularnewline
\bottomrule
\end{longtable}

启动器清除继承的实验变量。 最终服务器只接收模型、别名、主机和端口参数。
运行日志和服务器属性确认了实际默认值。

\hypertarget{ux5339ux914dux6027ux80fdux7ed3ux679c}{%
\section{5. 匹配性能结果}\label{ux5339ux914dux6027ux80fdux7ed3ux679c}}

性能请求使用固定的完整输入词元列表和 65 个固定输出词元。
单词元偏置使两个版本的解码历史完全相同。
这一合成负载只测性能，答案质量由单独的自然回答请求检查。

所有请求的提示缓存复用为 0。 测试不清理操作系统文件缓存。
每次请求前重新启动服务器，加载时间不计入请求计时。 28
次请求全部完成，每次都生成 65 个词元。
输入散列、输出词元、全文和计时都符合协议。

短输入每形状 3 对，两臂顺序交替。 8192 输入每个模型 1
对，两个模型采用相反顺序。
表中数值为每臂请求指标的算术均值，斜线前为起点，斜线后为最终版。
时间减少和吞吐增加使用不同的计算公式。

\begin{longtable}[]{@{}lrrrrr@{}}
\toprule
模型 & 输入 & 配对数 & 首词元时间，s & 服务端预填充，s &
预填充，词元/s\tabularnewline
\midrule
\endhead
2B & 512 & 3 & 23.835／19.556 & 23.831／19.551 &
21.485／26.188\tabularnewline
2B & 2048 & 3 & 115.639／81.633 & 115.632／81.626 &
17.712／25.090\tabularnewline
2B & 8192 & 1 & 846.321／408.162 & 846.303／408.145 &
9.680／20.071\tabularnewline
4B & 512 & 3 & 63.697／49.865 & 63.692／49.860 &
8.039／10.269\tabularnewline
4B & 2048 & 3 & 354.441／213.497 & 354.434／213.490 &
5.778／9.593\tabularnewline
4B & 8192 & 1 & 3234.172／1126.535 & 3234.155／1126.518 &
2.533／7.272\tabularnewline
\bottomrule
\end{longtable}

\begin{longtable}[]{@{}lrrrrr@{}}
\toprule
模型 & 输入 & 解码，词元/s & 吞吐变化 & 请求时长，s & 采样峰值
RSS，MiB\tabularnewline
\midrule
\endhead
2B & 512 & 3.387／4.292 & +26.69\% & 43.031／34.714 &
2471.97／2469.73\tabularnewline
2B & 2048 & 3.108／4.186 & +34.68\% & 136.554／97.169 &
2472.29／2469.96\tabularnewline
2B & 8192 & 2.239／3.600 & +60.82\% & 875.359／426.221 &
2473.30／2471.19\tabularnewline
4B & 512 & 1.315／2.133 & +62.17\% & 113.128／80.346 &
5342.20／5340.58\tabularnewline
4B & 2048 & 1.133／1.965 & +73.37\% & 411.793／246.579 &
5342.20／5340.88\tabularnewline
4B & 8192 & 0.724／1.485 & +104.91\% & 3323.894／1170.324 &
5342.88／5342.11\tabularnewline
\bottomrule
\end{longtable}

\begin{longtable}[]{@{}lrrrr@{}}
\toprule
模型 & 输入 & 首词元时间减少 & 预填充吞吐变化 &
解码耗时减少\tabularnewline
\midrule
\endhead
2B & 512 & 17.96\% & +21.89\% & 21.04\%\tabularnewline
2B & 2048 & 29.41\% & +41.66\% & 25.72\%\tabularnewline
2B & 8192 & 51.77\% & +107.35\% & 37.82\%\tabularnewline
4B & 512 & 21.72\% & +27.74\% & 38.34\%\tabularnewline
4B & 2048 & 39.77\% & +66.02\% & 42.32\%\tabularnewline
4B & 8192 & 65.17\% & +187.09\% & 51.20\%\tabularnewline
\bottomrule
\end{longtable}

RSS 指服务器进程的驻留内存。 表中为每次请求采样峰值的均值，采样间隔 0.2
s。 独立汇总保存了均值、中位数、最小值、最大值和每次请求的身份。
按冻结门槛，所有形状的首词元时间和解码耗时都没有明确退化。 门槛取 3\%
与起点臂完整波动范围中的较大值。

\hypertarget{ux6700ux7ec8ux7248ux5bb9ux91cfux6269ux5c55}{%
\subsection{5.1
最终版容量扩展}\label{ux6700ux7ec8ux7248ux5bb9ux91cfux6269ux5c55}}

单独的冻结协议把容量扩展到 20480 或 36864，对应实际输入 16384 或 32768
词元。 最终服务器除模型、别名、主机和端口外，只额外接收容量参数。
优化默认值不变，普通默认上下文仍为 12288。
两个模型的全部请求都完成，输出 65 个匹配词元，提示复用为 0。 16384
输入每个模型 2 对，32768 输入每个模型 1 对。

\begin{longtable}[]{@{}lrrrrr@{}}
\toprule
模型 & 输入 & 配对数 & 首词元时间，s & 预填充，s &
预填充，词元/s\tabularnewline
\midrule
\endhead
2B & 16384 & 2 & 2829.646／1034.907 & 2829.614／1034.876 &
5.797／15.832\tabularnewline
2B & 32768 & 1 & 10761.012／2983.902 & 10760.952／2983.845 &
3.045／10.982\tabularnewline
4B & 16384 & 2 & 11256.195／2975.716 & 11256.164／2975.685 &
1.456／5.506\tabularnewline
4B & 32768 & 1 & 42070.823／8866.022 & 42070.765／8865.966 &
0.779／3.696\tabularnewline
\bottomrule
\end{longtable}

\begin{longtable}[]{@{}lrrrr@{}}
\toprule
模型 & 输入 & 解码，词元/s & 请求时长，s & 峰值 RSS，MiB\tabularnewline
\midrule
\endhead
2B & 16384 & 1.659／2.914 & 2868.832／1057.226 &
2571.781／2569.617\tabularnewline
2B & 32768 & 1.051／2.141 & 10822.886／3014.282 &
2767.641／2766.426\tabularnewline
4B & 16384 & 0.490／1.117 & 11388.985／3033.905 &
5600.367／5600.611\tabularnewline
4B & 32768 & 0.291／0.752 & 42294.220／8952.505 &
6117.164／6117.809\tabularnewline
\bottomrule
\end{longtable}

\begin{longtable}[]{@{}lrrrr@{}}
\toprule
模型 & 输入 & 首词元时间减少 & 预填充提升 & 解码提升\tabularnewline
\midrule
\endhead
2B & 16384 & 63.43\% & +173.13\% & +75.59\%\tabularnewline
2B & 32768 & 72.27\% & +260.64\% & +103.72\%\tabularnewline
4B & 16384 & 73.56\% & +278.29\% & +128.23\%\tabularnewline
4B & 32768 & 78.93\% & +374.52\% & +158.34\%\tabularnewline
\bottomrule
\end{longtable}

独立检查核对了 12 次请求、221 份收据和 101704 个内存采样。 控制器总耗时
28.381 h，没有请求达到冻结的时间或内存上限。 RSS
采样覆盖服务器加载和请求，标称间隔 1 s。 采样到的最低系统可用内存为
11396.441 MiB，最高服务器 RSS 为 6117.809 MiB。
完整协议、样本范围和复现命令见\href{https://github.com/MoyaMryia/riscv-accl/blob/main/spacemit/reports/2026-10-16-ultralong-final-comparison.md}{容量报告}。

\hypertarget{ux6b63ux786eux6027ux4e0eux5931ux8d25ux8bb0ux5f55}{%
\section{6.
正确性与失败记录}\label{ux6b63ux786eux6027ux4e0eux5931ux8d25ux8bb0ux5f55}}

最终构建通过 104 项 route、95 项注意力、432 项 SSM 和 158 项 GDN
原生检查。 转储散列与已核验的 RC1 参考一致。 两个模型共 64
条全词表状态记录逐位一致。 这些记录覆盖 RS 0／3、状态清空，以及上下文
4096 下的 32／33／64／96 前缀。

6 次完整自然回答覆盖两个模型的起点、RC1 C0 和最终版三臂。
每次回答都自然停止，并通过冻结 Unicode 任务的 106 项功能检查。 2B
三臂全文一致，各生成 508 个词元。 4B 起点生成 595 个词元，全文不同。 RC1
与最终版的 4B 回答全文一致，各生成 521 个词元。 两种 4B
文本都通过功能检查。

r1 交付通过数值检查，但参数解析器的 GPU 默认值断言失败。 r2
修正了该断言，但启动时自动使用 4 个槽位。 r3
在源码中把服务器默认槽位固定为 1，并完整重新构建。
两个失败版本的源码、构建身份和日志都已保存。

旧的三项全开组合在 4B／2048／65 上相对 B1 退化
3.0220477\%，该结论在原测量条件下保留。 IME
标量调度候选未通过性能初筛，继续排除。 GPU
卸载和分页分配的内存收益没有合格结果。
预测器的误差和宽区间与实测模型结果分开记录。

最终搜索测试 B1-r2、C1 和 C2-r2 的 8 种组合，都没有启用 MTP。 此前的 MTP
试验在两个模型上都通过了相对质量门槛。 该试验的六题总耗时，2B 减少
8.63\%，4B 减少 15.22\%。 MTP
保持关闭，原因是该试验另设了绝对可用性门槛。 直接解码在区间代码题和 2B
算术题上同样未通过该门槛。
\href{https://github.com/MoyaMryia/riscv-accl/blob/main/spacemit/reports/2026-10-16-next-plan.md}{下一步计划}将在
r3 上按相对质量门槛和分类速度门槛重新测试 MTP。

\hypertarget{ux90e8ux7f72ux4e0eux8bc1ux636eux5305}{%
\section{7. 部署与证据包}\label{ux90e8ux7f72ux4e0eux8bc1ux636eux5305}}

成果包目录 \nolinkurl{final-release}
包含可运行代码、一键部署入口、官方转换脚本、部署说明、性能报告和设计说明。
它包含完整生产源码与二进制、固定的转换源码、模型身份元数据和独立证据归档。
生产二进制保持原核验散列。 新入口先创建 tmux
会话并取得共用锁，再启动服务器。
两个模型的启动都通过身份、实际默认值、锁和健康检查。

\begin{Shaded}
\begin{Highlighting}[]
\FunctionTok{bash}\NormalTok{ deploy.sh /absolute/path/to/model.gguf}
\end{Highlighting}
\end{Shaded}

成果包清单、检查表和文件散列记录了完整内容。
\href{https://github.com/MoyaMryia/riscv-accl/blob/main/spacemit/final-release/SUBMISSION-PLAN.md}{交付计划}列出六项材料的对应位置。

归档版本也可以直接启动。
两个模型使用同一个二进制和同一组默认值，只更换模型路径。 在 tmux
中，从解压后的版本目录执行：

\begin{Shaded}
\begin{Highlighting}[]
\FunctionTok{bash}\NormalTok{ package/with{-}lock.sh python3 package/start.py {-}{-}model /path/to/model.gguf}
\end{Highlighting}
\end{Shaded}

依赖目录不同时，使用 \nolinkurl{--spert} 和 \nolinkurl{--tcm}
指定库位置。 这两个参数只选择库位置，不选择性能配置。
启动器核对二进制、模型和运行库散列。
二进制迁移后，两个模型的启动都通过检查。

归档为
\nolinkurl{lab/final-default-20261013/release/k1-final-20261013-r3.tar.gz}。
外部旁注文件记录归档 SHA-256。 Git
文本视图省略大型转储、线程采样和二进制。
严格重放收据需要保留的完整归档。 本地检查核对 338 份对照收据和 119
份最终验证收据。 搜索检查另核对 951 份搜索收据和 19 份原生收据。

\hypertarget{ux5386ux53f2ux8bc1ux636e}{%
\section{8. 历史证据}\label{ux5386ux53f2ux8bc1ux636e}}

\href{https://github.com/MoyaMryia/riscv-accl/blob/main/spacemit/reports/snapshots/2026-10-12/SUBMISSION-REPORT.pdf}{10
月 12 日快照}保存旧版报告的全部内容。
不同对照下的历史百分比不能相加或相乘。
\href{https://github.com/MoyaMryia/riscv-accl/blob/main/spacemit/reports/2026-10-07-release-verification.md}{10
月 7 日发布}和
\href{https://github.com/MoyaMryia/riscv-accl/blob/main/spacemit/reports/2026-10-13-unified-release-candidate.md}{10
月 13 日 RC1} 保留为对照记录。
\href{https://github.com/MoyaMryia/riscv-accl/blob/main/spacemit/reports/WORK-LEDGER.md}{项目工作总账}列出全部工作、采用项和每个停止事项的原因。
\href{https://github.com/MoyaMryia/riscv-accl/blob/main/spacemit/reports/2026-10-16-next-plan.md}{下一步计划}按预期收益排列剩余实验。

\hypertarget{ux540eux7eedux9a8cux6536ux8bb0ux5f55}{%
\section{9. 后续验收记录}\label{ux540eux7eedux9a8cux6536ux8bb0ux5f55}}

下列报告记录 2026-10-16 之后的独立验收。第 5 节与 5.1
节的冻结测量保持不变。

\begin{itemize}
\tightlist
\item
  官方检查点加工：两个检查点完成转换、量化和打包；张量、类型与分词器数据与交付文件一致，对话模板不同（\href{https://github.com/MoyaMryia/riscv-accl/blob/main/spacemit/reports/2026-10-21-official-conversion.md}{报告}）。
\item
  固定多轮 v2：48 轮请求全部完成；独立复核 2B 15/24、4B 19/24
  正确，用户接受测得质量并保留已知错误（\href{https://github.com/MoyaMryia/riscv-accl/blob/main/spacemit/reports/2026-10-21-fixed-multi-turn-v2-results.md}{报告}）。
\item
  固定多轮 v1：48 个计划轮次完成 32 个，退出码 2；6 个会话通过、6
  个失败（\href{https://github.com/MoyaMryia/riscv-accl/blob/main/spacemit/reports/2026-10-21-fixed-multi-turn-results.md}{报告}）。
\item
  自然长上下文：短、16k、32k 档位共 12
  次回答全部正确（\href{https://github.com/MoyaMryia/riscv-accl/blob/main/spacemit/reports/2026-10-22-long-context-natural-results.md}{报告}）。
\item
  部分请求 MTP：文档字符串规则在代码网格上净收益 3.48\%，族区间
  1.30\%～5.10\%（\href{https://github.com/MoyaMryia/riscv-accl/blob/main/spacemit/reports/2026-10-21-2b-mtp-docstring-rule-validation.md}{报告}）。
\end{itemize}

\hypertarget{ux9002ux7528ux8303ux56f4}{%
\section{适用范围}\label{ux9002ux7528ux8303ux56f4}}

配置：MUSE-Pi-Pro，K1／X60，标称 16 GB 内存，版本
\nolinkurl{k1-final-20261013-r3}，4 线程，批大小 32，F16
键值缓存，单槽。 样本：512 和 2048 输入每形状 3 对；8192 和 32768
输入每个模型 1 对；16384 输入每个模型 2 对。
自然回答检查每个模型每臂使用一个 Unicode 代码任务。
默认配置是冻结搜索范围内的均衡选择。
整体对照没有分离各项已采用路径的贡献。

\begin{itemize}
\tightlist
\item
  待验证：8192、16384 和 32768 输入的重复稳定性。阻塞：这些形状只有 1～2
  对。解除条件：每形状至少测量 3 对。
\item
  待验证：长状态正确性。阻塞：r3 上没有长状态检查。解除条件：在 r3
  上执行长状态检查，并与自然长文回答分开记录。
\item
  待验证：宽 RVV、route 3、B1-r2 和 C1
  各自的贡献。阻塞：最终对照只比较整体版本。解除条件：在同一 r3
  二进制上逐项关闭并配对测量。
\item
  待验证：确切依赖的统一公开下载归档。阻塞：厂商运行库不在交付包中。解除条件：发布确切运行库文件及其
  SHA-256。
\item
  待验证：8 GB 设备上的运行。阻塞：没有测试 8 GB 板。解除条件：在 8 GB
  板上执行容量协议。
\end{itemize}
\end{document}
